% Part: propositional-logic

\documentclass[../../include/open-logic-part]{subfiles}

\begin{document}

\olpart{pl}{Logika Proposisional}

\begin{editorial}
  Bagean iki ngemot materi ngenani logika proposisional klasik. Bab
  kapisan rada dhasar lan mung ndhaptar definisi lan asil;
  akeh bukti ora ditindakake, nanging dipasrahake dadi gladhen.
  Materi ngenani sistem bukti lan teorema kelengkapan dijupuk saka
  bagean logika tataran kapisan, kanthi tag ``FOL'' dipateni.
  Mula kabeh kang gegayutan karo predikat, term, lan kuantor ora
  dilebokake, lan pirembugan ngenani !!{structure}s~$\Struct{M}$ diganti
  dadi pirembugan ngenani !!{valuation}s~$\pAssign{v}$.

  Ana rancangan kanggo ngrembakakake bagean iki kanthi rerincen kang
  luwih akeh, lan nambah topik sarta asil liyane, kayata kelengkapan
  fungsional tumrap kabeneran.
\end{editorial}

\olimport[syntax-and-semantics]{syntax-and-semantics}

\tagfalse{FOL}

\olimport[../first-order-logic/proof-systems]{proof-systems}

\iftag{prfSC}{%
  \olimport[../first-order-logic/sequent-calculus]{sequent-calculus}
}{}

\iftag{prfND}{%
  \olimport[../first-order-logic/natural-deduction]{natural-deduction}
}{}

\iftag{prfTab}{%
  \olimport[../first-order-logic/tableaux]{tableaux}
}{}

\iftag{prfAX}{%
  \olimport[../first-order-logic/axiomatic-deduction]{axiomatic-deduction}
}{}

\olimport[../first-order-logic/completeness]{completeness}

\tagtrue{FOL}

\OLEndPartHook

\end{document}
